% TOP_PROBLEM: {"id": 30006699, "title": "Permanent versus Determinant over C", "category_id": 15, "category": {"id": 15, "name": "computer_science", "display_name": "Computer Science", "description": "Computational complexity, algorithms, and theoretical CS.", "slug": "computer-science", "order_index": 15, "created_at": "2026-07-31T15:26:25.671Z"}, "status": "open"}
% FIRST_POSTED: 2026-09-27
% !TeX program = lualatex
% Compile twice: lualatex 30006699.tex
% Original entry retained; research subsection and [E] references added.
% No external bibliography, image, data, or style file is required.
\documentclass[11pt,a4paper]{article}
\usepackage[margin=24mm,headheight=14pt]{geometry}
\usepackage{fontspec}
\setmainfont{Latin Modern Roman}
\setsansfont{Latin Modern Sans}
\setmonofont{Latin Modern Mono}
\usepackage{amsmath,amssymb,mathtools,amsthm}
\usepackage{unicode-math}
\setmathfont{Latin Modern Math}
\usepackage{microtype}
\usepackage{enumitem,xurl,needspace}
\usepackage[dvipsnames]{xcolor}
\usepackage{hyperref}
\hypersetup{unicode=true,colorlinks=true,linkcolor=MidnightBlue,
 urlcolor=MidnightBlue,bookmarksnumbered=true,
 pdftitle={Research notebook - Problem 30006699},
 pdfauthor={Research attempt; not independently refereed}}
\usepackage{bookmark}
\usepackage{titlesec}
\titleformat{\section}{\Large\bfseries\raggedright}{\thesection}{0.7em}{}
\usepackage{fancyhdr}
\pagestyle{fancy}\fancyhf{}
\fancyhead[L]{\small\sffamily Research notebook}
\fancyhead[R]{\small\sffamily Problem 30006699}
\fancyfoot[L]{\scriptsize 27 September 2026}
\fancyfoot[C]{\thepage}
\fancyfoot[R]{\scriptsize Unrefereed research attempt}
\setlength{\parindent}{0pt}
\setlength{\parskip}{5pt plus 1pt}
\setlength{\emergencystretch}{3em}
\clubpenalty=10000
\widowpenalty=10000
\displaywidowpenalty=10000
\setcounter{secnumdepth}{1}
\setlist[itemize]{leftmargin=3em,itemsep=5pt,topsep=4pt,parsep=0pt}
\allowdisplaybreaks
\newcommand{\N}{\mathbb{N}}
\newcommand{\Z}{\mathbb{Z}}
\newcommand{\Q}{\mathbb{Q}}
\newcommand{\R}{\mathbb{R}}
\newcommand{\C}{\mathbb{C}}
\newcommand{\F}{\mathbb{F}}
\newcommand{\A}{\mathbb{A}}
\newcommand{\PP}{\mathbb P}
\newcommand{\E}{\mathbb{E}}
\newcommand{\eps}{\varepsilon}
\newcommand{\dd}{\,\mathrm d}
\newcommand{\sref}[2]{\hyperlink{source:#1:#2}{[S#2]}}
\newcommand{\eref}[2]{\hyperlink{extra:#1:#2}{[E#2]}}
\newif\ifshowcatalogue
\showcataloguetrue
\newcommand{\cataloguescope}[1]{\ifshowcatalogue
 \paragraph{Exact scope in the supplied catalogue.}
 \begin{quote}\small #1\end{quote}\fi}
\newtheorem{notebooklemma}{Lemma}[section]
\newtheorem{notebookproposition}[notebooklemma]{Proposition}
\numberwithin{equation}{section}
\begin{document}
\setcounter{section}{0}
% ================================================================
% problemId: 30006699
% Canonical problem ID: 30006699
% exactTarget SHA-256: e7d97b19c601e89a6a4dc884103b4b436ef184b0bc38158b2e6c93e31427bf88
\clearpage
\section*{\texorpdfstring{Permanent versus Determinant over C}{Permanent versus Determinant over C}}\label{30006699}
\subsection{Definitions and mathematical statement}
For the permanent $\operatorname{per}_n(X)=\sum_{\sigma\in S_n}\prod_i x_{i,\sigma(i)}$, define determinantal complexity by
\[
 \operatorname{dc}(\operatorname{per}_n)=\min\{m:\operatorname{per}_n(X)=\det A(X),\
 A(X)\in\operatorname{Mat}_m(\C[x]),\ A_{ij}\text{ affine-linear}\}.
\]
The assertion is $\forall c,C>0\ \forall N\ \exists n\ge N:\operatorname{dc}(\operatorname{per}_n)>Cn^c$. The identity must hold exactly as a polynomial. Approximation by determinants, symmetry-restricted representations, and general arithmetic-circuit complexity are different measures. The chosen complexity-class formulation is $\mathsf{VP}_{\rm ws}\ne\mathsf{VNP}$, not an asserted equivalence with $\mathsf{VP}\ne\mathsf{VNP}$.
\par\smallskip\textit{Formulation and concept references:} \sref{30006699}{1}.
\cataloguescope{Over C, let dc(per\_n) be the least m for which the n\ensuremath{\times}n permanent per\_n(X) equals the determinant of an m\ensuremath{\times}m matrix whose entries are affine-linear forms in the n\textasciicircum{}2 variables X, as an exact polynomial identity. Assert that dc(per\_n) is not O(n\textasciicircum{}c) for any real constant c>0: explicitly, for every c>0, C>0 and integer N\ensuremath{\ge}1, some n\ensuremath{\ge}N satisfies dc(per\_n)>C*n\textasciicircum{}c. No restrictions of equivariance, sparsity or symmetry are imposed on the representing matrix. This is exact, not border, determinantal complexity over C; the selected assertion is the source’s VP\_ws != VNP formulation, not an assertion that it is equivalent to VP != VNP.}
\subsection{Short English statement}
Must every exact determinant representation of the permanent eventually exceed every polynomial size bound?
\subsection{Research attempt}

\paragraph{Current frontier.}
The general characteristic-zero lower bound supplied by the Mignon--Ressayre Hessian method is quadratic, while known unrestricted upper bounds are exponential; Bogart--Valero review the bounds and extend the technique to other permanent-like polynomials \eref{30006699}{1}. Exponential lower bounds obtained with symmetry requirements do not apply to the unrestricted affine-linear representations in this entry \eref{30006699}{2}. There is also a precise barrier to a proposed geometric-complexity route: B\"urgisser--Ikenmeyer--Panova, Theorem 1.4, rule out the relevant occurrence obstructions when the determinant size is at least the twenty-fifth power of the permanent size. Their statement does not rule out all multiplicity comparisons \eref{30006699}{3}. The June 2026 conormal-specialization manuscript in the catalogue concerns a different polynomial and border complexity; its claimed bound is not an unrestricted permanent lower bound \eref{30006699}{4}.

\paragraph{Attack route.}
A single Hessian cannot exceed rank $n^2$, so repeating the rank estimate at one point cannot produce a superpolynomial bound. Representation-theoretic multiplicity obstructions remain conceivable, but an occurrence-only argument would meet the cited barrier. The route investigated here is to force \emph{all higher jets to share one affine matrix pencil}. The additional compatibility might cost more than any individual derivative tensor. We derive the exact constraints, test a noncommutative shortcut, and prove why keeping the jet order fixed cannot suffice.

\paragraph{Attempt: an exact smooth point and its full Hessian.}
Let $f=\operatorname{per}_n$, initially $n\geq3$, and let $P$ have every entry equal to $1$ except $P_{nn}=1-n$. Then
\[
 f(P)=n!-n(n-1)!=0,
 \qquad \frac{\partial f}{\partial x_{nn}}(P)=(n-1)!\ne0.
\]
Index the Hessian $H$ by pairs $(i,j)$. Its entries are permanents of complementary minors, so
\begin{equation}\label{r57:hessian}
 H_{ij,kl}=\begin{cases}
 0,&i=k\text{ or }j=l,\\
 -2(n-3)!,&i\ne k,\ j\ne l,\ n\notin\{i,k,j,l\},\\
 (n-2)!,&\text{otherwise, when }i\ne k,\ j\ne l.
 \end{cases}
\end{equation}
For the second case the remaining minor is all ones with a single entry reduced by $n$, and its permanent is $(n-2)!-n(n-3)!$. Thus the signs in this formula do not come from a numerical fit.

\begin{notebooklemma}\label{r57:fullrank}
The Hessian in \eqref{r57:hessian} has rank $n^2$.
\end{notebooklemma}
\begin{proof}
Divide the matrix by $(n-3)!$, and write a putative kernel vector as an $(n-1)\times(n-1)$ block $A$, last-row entries $b_j$, last-column entries $c_i$, and corner entry $u$. Put $m=n-1$ and $h=n-2>0$. The equation in the corner row gives $\sum_{ij}A_{ij}=0$. The equations in the other last-row positions give
\[
 \sum_i A_{ij}=\sum_i c_i\quad\text{for every }j.
\]
Summing over $j$ forces $\sum_i c_i=0$, and then every column sum of $A$ is zero. The last-column equations similarly give zero row sums and $\sum_jb_j=0$. The remaining equations reduce to
\[
 -2A_{ij}+h(u-b_j-c_i)=0.
\]
Summing one of these rows shows $c_i=u$ for every $i$. Since the $c_i$ sum to zero, $u=0$ and every $c_i=0$. Column sums similarly give every $b_j=0$, and finally $A=0$.
\end{proof}

For $n=2$ the Hessian is directly a nonsingular $4\times4$ permutation matrix. The following argument consequently also covers that edge case. This recovers the known quadratic mechanism; it is not being presented as a new quadratic lower bound \eref{30006699}{1}.

Suppose an exact representation $f(X)=\det A(X)$ has size $M$. At the smooth zero $P$, the matrix $A(P)$ has rank exactly $M-1$: rank at most $M-2$ would make its adjugate zero, hence make every first derivative of the pullback determinant zero. Constant invertible row and column operations yield a nonzero scalar $\lambda$ and linear forms arranged as
\begin{equation}\label{r57:normalform}
 f(P+z)=\lambda\det\begin{pmatrix}
 a(z)&b(z)^t\\ c(z)&I_r+D(z)
 \end{pmatrix},\qquad r=M-1.
\end{equation}
Here $a$, all entries of $b,c$, and all entries of $D$ are homogeneous linear forms in the $n^2$ displacement variables $z$. In particular $a=df_P/\lambda$.

The quadratic part is $\lambda(a\operatorname{tr}D-b^tc)$, so its Hessian has rank at most $2+2r=2M$. Lemma~\ref{r57:fullrank} therefore yields
\begin{equation}\label{r57:quadraticbound}
 \operatorname{dc}(\operatorname{per}_n)\geq\lceil n^2/2\rceil.
\end{equation}
Restricting to the tangent hyperplane $T=\ker df_P$ does not improve this. Homogeneity gives $HP=(n-1)df_P$, while $df_P(P)=nf(P)=0$. Thus the radical of the restricted Hessian on $T$ is exactly $\C P$, and its rank is $n^2-2$. The available upper bound there is $2r$, giving the same inequality.

\paragraph{Attempt: compatibility of all higher orders.}
The full, exact polynomial identity behind \eqref{r57:normalform} is
\begin{equation}\label{r57:schur}
 f(P+z)=\lambda\bigl(a\det(I_r+D)-b^t\operatorname{adj}(I_r+D)c\bigr).
\end{equation}
It follows either by a block determinant calculation over the rational function field, or by multiplying the Schur complement by the determinant and then using polynomial identity. Restricting to $T$, where $a=0$, gives a formal power series with constant denominator $1$:
\begin{equation}\label{r57:series}
 -\frac{f(P+z)|_T}{\lambda\det(I_r+D|_T)}
   =b^t(I_r+D)^{-1}c\bigm|_T
   =\sum_{k\geq0}(-1)^k b^tD^kc\bigm|_T.
\end{equation}
Each coefficient of a fixed degree involves only finitely many terms, so this requires no analytic convergence or exchange of limiting operations.

All homogeneous pieces on the right use \emph{the same} $r\times r$ matrix $D$ and the same two vectors. Write $v_k=b^tD^kc$, and let $e_j(D)$ be the coefficients of the characteristic polynomial of $D$. Cayley--Hamilton supplies
\begin{equation}\label{r57:recurrence}
 v_{k+r}-e_1(D)v_{k+r-1}+e_2(D)v_{k+r-2}
          -\cdots+(-1)^re_r(D)v_k=0\qquad(k\geq0).
\end{equation}
Every summand has degree $k+r+2$. These are compatibility identities in the commutative polynomial ring, not independent rank estimates that may be added together. Also the denominator in \eqref{r57:series} depends on the unknown representation: discarding it would change the invariants being compared.

To make a potential lower-bound statement precise, define $J_t(f,P)$ to be the least $r$ for which linear data $a,b,c,D$ and $\lambda\ne0$ satisfy \eqref{r57:schur} through Taylor order $t$ at $P$, with the displayed block sizes. This minimum is over \emph{all} such data; no symmetry or sparsity is required. Every exact representation gives
\begin{equation}\label{r57:jetbound}
 J_t(f,P)+1\leq\operatorname{dc}(f).
\end{equation}
Thus a proof that $J_{t(n)}(\operatorname{per}_n,P_n)$ exceeds every polynomial bound, with the same quantifier meaning as in the original statement, would prove the target. Equations~\eqref{r57:series}--\eqref{r57:recurrence} specify constraints for such an attack. No superpolynomial lower bound for this minimum is established here.

There is a further limitation that can be proved rather than guessed.
\begin{notebooklemma}[Fixed jet orders have a polynomial ceiling]\label{r57:fixedorder}
For fixed $t\geq1$, a polynomial germ in $v$ variables with zero constant term and nonzero differential has
\[
 J_t\leq O\!\left(t\binom{v+t}{t}\right).
\]
In particular, no fixed-order version of \eqref{r57:jetbound} can yield a superpolynomial lower bound in $n$ when $v=n^2$.
\end{notebooklemma}
\begin{proof}
Its Taylor polynomial $q$ of degree at most $t$ has at most $\binom{v+t}{t}$ monomials. Represent each nonzero monomial by a separate directed path, with a common source and sink, multiplying its coefficient into one edge. The resulting acyclic graph has $O(t\binom{v+t}{t})$ vertices and affine-linear edge labels, and its path sum is $q$.

If $B$ is its strictly upper triangular adjacency matrix, then $I-B$ has determinant $1$ and
$e_s^t(I-B)^{-1}e_t=q$. Therefore the affine-linear bordered matrix
\[
 \begin{pmatrix}0&e_s^t\\-e_t&I-B\end{pmatrix}
\]
has determinant $q$. Since $q(0)=0$ and $dq_0\ne0$ for the germ under consideration, this representation has corank one at zero and can be put into \eqref{r57:normalform}. It reproduces the required jet. Constants in the bound are immaterial for its polynomial consequence.
\end{proof}
The zero constant term is part of the germ here because $P$ is on the hypersurface; that qualification is needed in applying the last normalization. The lemma forces $t(n)$ to become unbounded, or forces an attack using information beyond any one fixed jet order.

\paragraph{Stress test: a false noncommutative shortcut.}
A tempting idea is to put each monomial in row order, interpret the variables as noncommuting, and use the large ranks of word-coefficient matrices as a lower bound on $r$. This cannot be done without proving that the lift preserves low commutative determinantal complexity.

Here is an exact countertest. For the row-ordered lift of the permanent, form the matrix whose rows are length-$k$ prefixes and whose columns are length-$(n-k)$ suffixes. The entry is the coefficient of their concatenation. Prefixes using a column set $S$ can pair only with suffixes using its complement. There is one nonzero block for each of the $\binom nk$ sets $S$, and each block is an all-ones matrix. The rank is therefore $\binom nk$.

For the row-ordered lift of the determinant the rank is \emph{also} $\binom nk$. In the same block, the permutation sign factors as the sign of the prefix ordering, the sign of the suffix ordering, and the fixed shuffle sign determined by $S$. Each block is again a nonzero outer product of two vectors and has rank one. Thus this lift has exponential middle-cut rank even though the original commutative determinant has its tautological size-$n$ representation. Any claimed transfer assigning that rank as a mandatory cost to commutative determinant size is false. This is a counterexample to the shortcut, not to the permanent conjecture.

The accompanying exact-arithmetic script checks the Hessian ranks for $n=2,\ldots,7$, obtaining $4,9,16,25,36,49$. It computes both ordered-lift ranks for $n=2,3,4,5$ at $k=\lfloor n/2\rfloor$, obtaining respectively $2,3,6,10$ for each lift. It also verifies a symbolic Schur-adjugate identity and a $2\times2$ Cayley--Hamilton instance. These computations test signs and conventions; the preceding proofs supply the general claims.

\paragraph{Outcome.}
\textbf{Outcome: Exploratory approach with unresolved gap.}

The notebook derives exact shared-pencil constraints on higher jets, proves a polynomial ceiling for fixed-order jet tests, and exhibits the determinant itself as a countertest to the row-order lifting shortcut. The only lower bound on the target obtained here is the already known quadratic bound. The unresolved step is a superpolynomial lower bound for compatible unbounded-order data, allowing every affine-linear representation and its representation-dependent determinant unit. No border statement, symmetry-restricted result, or arithmetic-circuit separation is substituted for the requested exact $\mathsf{VP}_{\rm ws}\ne\mathsf{VNP}$ statement.

\subsection{Sources}
\begin{itemize}\raggedright
\item[\textnormal{[S1]}]\hypertarget{source:30006699:1}{} Bürgisser, Ikenmeyer and Panova, arXiv:1604.06431v3, Conjecture 1.1; Landsberg and Ressayre, arXiv:1508.05788v2, Conjecture 1.1 and affine-linear representation definition.
\url{https://arxiv.org/pdf/1604.06431}.
{\small\textit{Catalogue formulation source.}}
\item[\textnormal{[S2]}]\hypertarget{source:30006699:2}{} Bounds on Determinantal Complexity of Two Types of Generalized Permanents.
\url{https://arxiv.org/html/2202.13016}.
{\small\textit{Catalogue research/status reference; not a proof endorsement.}}
\item[\textnormal{[S3]}]\hypertarget{source:30006699:3}{} A near-quadratic lower bound on the border determinantal complexity of sum\_i x\_i\textasciicircum{}n via conormal specialization.
\url{https://arxiv.org/abs/2606.13628}.
{\small\textit{Catalogue research/status reference; not a proof endorsement.}}

% Research references added for this notebook.
\item[\textnormal{[E1]}]\hypertarget{extra:30006699:1}{} Tristram Bogart and Juan Andr\'es Valero, \emph{Bounds on Determinantal Complexity of Two Types of Generalized Permanents}, arXiv:2202.13016, Introduction and Section 3; inspected HTML has internal date 24 August 2026.
\url{https://arxiv.org/html/2202.13016}.
{\small\textit{Review of the Mignon--Ressayre quadratic bound, Grenet upper bound, and Hessian method; the notebook does not claim a new quadratic lower bound. Consulted 27 September 2026.}}
\item[\textnormal{[E2]}]\hypertarget{extra:30006699:2}{} J. M. Landsberg and Nicolas Ressayre, \emph{Permanent v. determinant: an exponential lower bound assuming symmetry and a potential path towards Valiant's conjecture}, arXiv:1508.05788v2, Conjecture 1.1, definitions, and symmetry-restricted results.
\url{https://arxiv.org/html/1508.05788v2}.
{\small\textit{The symmetry qualification is not removed in applying or interpreting this source. Consulted 27 September 2026.}}
\item[\textnormal{[E3]}]\hypertarget{extra:30006699:3}{} Peter B\"urgisser, Christian Ikenmeyer and Greta Panova, \emph{No occurrence obstructions in geometric complexity theory}, arXiv:1604.06431v3 (17 September 2018), Theorem 1.4 and the following discussion of multiplicity obstructions.
\url{https://arxiv.org/html/1604.06431v3}.
{\small\textit{Exact scope of the occurrence-obstruction barrier; it does not rule out every representation-theoretic method. Consulted 27 September 2026.}}
\item[\textnormal{[E4]}]\hypertarget{extra:30006699:4}{} Karthik Sheshadri, \emph{A near-quadratic lower bound on the border determinantal complexity of $\sum_i x_i^n$ via conormal specialization}, arXiv:2606.13628v1 (11 June 2026), abstract and main claims.
\url{https://arxiv.org/abs/2606.13628v1}.
{\small\textit{Different polynomial family and border model. This claimed result is not an input to any proof in the notebook and has not been independently verified here. Consulted 27 September 2026.}}
\end{itemize}

\end{document}
